\documentclass{tietreport}

% Import images
\usepackage{graphicx}
\graphicspath{{figures/}}

% Tables and maths
\usepackage{amsmath}
\usepackage{array}
\usepackage{tabularx}

% Bibliography configuration
\usepackage[
  date=year,
  style=alphabetic,
  backend=biber,
  sorting=nty,
  maxnames=5,
  minnames=3,
  maxalphanames=4,
  minalphanames=3,
  backref=true,
  doi=false,
  isbn=false,
  url=true,
  eprint=false
]{biblatex}
\addbibresource{references.bib}

%% ==================================================
%% PROJECT DETAILS
%% ==================================================

\TitlePageHeader{%
  UCS503P%
}

\ProjectTitle{%
  Personal Finance Manager%
}

\ProjectSubTitle{%
  A Desktop-Based Financial Tracking and Budgeting System%
}

\author{%
  \texttt{1024030022} Arjun Puri \\
  \texttt{1024030105} Anshdeep Singh%
}

\TitlePageSubText{%
  Group: \texttt{3C11} BE Third Year, CSE%
}

\AdvisorName{%
  Ms. Parmveer%
}

\TitlePageFooterText{{%
    \large Mid-Semester Evaluation \\[0.2cm]
    \bfseries Computer Science and Engineering
    Department%
  } \\[0.15cm]
  Thapar Institute of Engineering and Technology,
  Patiala%
}

%% ==================================================

\begin{document}

% Removes the blank pages caused by chapters opening on odd pages.
% Delete this line if your department requires chapters on right-hand pages.
\let\cleardoublepage\clearpage

\maketitle

\tableofcontents
\listoffigures
\listoftables

% ===== CHAPTER 1: INTRODUCTION =====

\chapter{Introduction}

\section{Background and Context}

Students and young professionals usually earn from several
sources, such as an allowance, a stipend or part-time work, and
spend across many categories, including food, travel,
subscriptions, shopping and bills. Without a structured system it
is hard to see where money is going or whether monthly spending is
under control. The \emph{Personal Finance Manager} (PFM) is a
desktop application that brings income, expenses, budgets and
savings goals into one place and turns them into summaries that
the user can read at a glance.

The project was proposed in~\cite{pfmproposal} with a
\emph{time-to-value} approach: build the login and
transaction-and-database core first, validate it, and then layer
budgets, goals and reports on top. This report documents the
implemented desktop application, its design, and the results of
the first phase.

\subsection{Problem Statement}

Tracking personal finances by hand or across scattered sources
leads to four recurring problems:

\begin{itemize}
  \item \textbf{Fragmentation:} transactions are spread across
    bank alerts, UPI apps and cash, with no single source of truth.
  \item \textbf{Low visibility:} there is no immediate view of how
    actual spending compares with a monthly budget.
  \item \textbf{No goal tracking:} savings targets, such as a
    laptop or a phone, are rarely tracked with a visible progress
    measure.
  \item \textbf{Manual, repetitive entry:} recurring costs such as
    rent are re-entered every month.
\end{itemize}

In its current version PFM provides a single MySQL-backed store for
records the user enters, which partly addresses fragmentation; it
does not import data from banks or UPI apps. It addresses low
visibility and goal tracking directly. Manual entry is not
reduced: recurring-expense handling is planned for a later
iteration (Section~\ref{sec:limitations}).

\section{Project Objectives}

\begin{enumerate}
  \item Provide a login and a single MySQL-backed store for a
    user's income and expense records, with validation of amount,
    date and required fields.
  \item Show the user's financial position (income, expenses, net
    savings, category breakdown and recent activity) on a
    dashboard that updates from the stored data.
  \item Let the user set monthly spending limits per category and
    see usage against each limit.
  \item Let the user create savings goals with a target, amount
    saved and optional deadline, and view progress towards each
    goal.
  \item Provide monthly reports with a category breakdown and a
    CSV export.
  \item Keep the presentation, logic and data layers separate so
    the system can later move to a web back end without
    redesigning the schema.
\end{enumerate}

Objectives~1--5 are implemented in the current build, although the
login does not yet hash passwords (Section~\ref{sec:limitations}).
Objective~6 is partly met and is discussed under future work.

% ===== CHAPTER 2: METHODOLOGY & DESIGN =====

\chapter{Methodology and Design}

\section{System Architecture}

\subsection{High-Level Design}

PFM is designed around the three-tier structure set out in the
proposal. The tiers are only partly separated in the current
build (Section~\ref{sec:limitations}).

\begin{itemize}
  \item \textbf{Presentation tier:} a Tkinter GUI with a login
    screen and eight pages reachable from a persistent sidebar
    (Dashboard, Transactions, Income, Expenses, Budget, Savings
    Goals, Reports and Settings).
  \item \textbf{Business-logic tier:} Python code that validates
    input and computes the totals, percentages and budget figures
    shown on screen. Some of this code still lives inside the GUI
    classes.
  \item \textbf{Data tier:} a MySQL database accessed through
    PyMySQL, with queries parameterised using placeholders.
\end{itemize}

The core workflow is: the user logs in; the dashboard loads;
the user adds an income or expense record; the input is
validated; the record is written to MySQL; and the totals,
budget usage and charts are recalculated when the views refresh.

\subsection{Navigation Architecture}

The application uses a single controller window (the \texttt{App}
class, derived from \texttt{tk.Tk}) that owns all window-level
settings. Each page is a \texttt{tk.Frame} built once and stacked
in a shared container with \texttt{grid()}. Switching pages calls
\texttt{show\_frame()}, which raises the chosen frame with
\texttt{tkraise()}, and each page exposes an \texttt{on\_show()}
hook that reloads its data when the page becomes visible. This
replaced an earlier design that destroyed and recreated widgets on
every navigation, which was slower and harder to maintain.

\subsection{Technical Stack}

\begin{itemize}
  \item \textbf{Language:} Python 3, using the standard library
    wherever possible.
  \item \textbf{GUI:} Tkinter~\cite{pythontkinter}, with the
    \texttt{tkcalendar} \texttt{DateEntry} widget for date
    selection~\cite{tkcalendardocs}.
  \item \textbf{Database:} MySQL~\cite{mysqlmanual} accessed
    through PyMySQL~\cite{pymysqldocs}.
  \item \textbf{Charts:} drawn directly on \texttt{tk.Canvas}
    (bar, line, donut and progress-bar views), with no external
    plotting library.
  \item \textbf{Export:} CSV export of the monthly report.
\end{itemize}

\subsection{Financial Calculations}

All displayed figures are deterministic formulas, so every number
can be audited by hand:

\begin{align*}
  \text{Net Savings} &= \text{Total Income} - \text{Total Expenses}\\
  \text{Budget Usage (\%)} &= \frac{\text{Amount Spent}}{\text{Budget Limit}} \times 100\\
  \text{Goal Progress (\%)} &= \frac{\text{Amount Saved}}{\text{Target Amount}} \times 100\\
  \text{Category Share (\%)} &= \frac{\text{Category Expense}}{\text{Total Expenses}} \times 100
\end{align*}

``Amount Spent'' is not defined identically everywhere. On the
Budget page it is the expense total over categories that have a
budget; on the sidebar spending-limit card it is the total of all
expenses. Percentages are shown as whole numbers, and the rounding
is not yet identical on every page. Both points are discussed in
Section~\ref{sec:analysis}.

% TODO: add the database design here (ER diagram and a table listing
% every table with its columns, types and keys). Send me the output of
% SHOW CREATE TABLE for each table and I will write this subsection.

\section{Implementation Details}

\subsection{Authentication}

The login screen (Figure~\ref{fig:login}) takes a username and
password, checks both fields are filled, and looks the user up
with a parameterised query on the \texttt{users} table. In the
current build the password is compared as plain text
(Section~\ref{sec:limitations}). A failed check shows an inline
error message; a database failure shows an error dialog. Pressing
Enter in the password field submits the form.

\begin{figure}[htbp]
  \centering
  \includegraphics[width=0.75\textwidth]{login.png}
  \caption{Login screen}
  \label{fig:login}
\end{figure}

\subsection{Dashboard}

The dashboard (Figure~\ref{fig:dashboard}) is the landing page. It
shows income, expenses and net savings with their change against
the previous month, a monthly overview chart for January to
December, an expense trend line, a list of recent transactions
with All / Income / Expenses filters, an expense-breakdown donut
chart and savings-goal progress. The sidebar also carries a
spending-limit card that shows the total spent against the monthly
limit. All charts are drawn on canvases and redraw when the
Refresh button is pressed.

\begin{figure}[htbp]
  \centering
  \includegraphics[width=\textwidth]{dashboard.png}
  \caption{Dashboard with summary statistics, charts, recent transactions and savings goals}
  \label{fig:dashboard}
\end{figure}

\subsection{Transactions}

The Transactions page (Figure~\ref{fig:transactions}) lists every
income and expense record in one scrollable table with date,
description, category, type and amount. Income rows are shown in
green and expense rows in red, and a footer reports the number of
records shown.

\begin{figure}[htbp]
  \centering
  \includegraphics[width=\textwidth]{transactions.png}
  \caption{Transactions page with colour-coded income and expense rows}
  \label{fig:transactions}
\end{figure}

\subsection{Income Management}

The Income page (Figure~\ref{fig:income}) shows all-time income and
the current month's income above a form with amount, category,
description, date and income source. Amount must be a number
greater than zero and the date must be a valid \texttt{dd-mm-yyyy}
value; invalid input triggers a warning and nothing is saved. The
category list covers salary, freelancing, business, investment,
interest and other common sources.

\begin{figure}[htbp]
  \centering
  \includegraphics[width=\textwidth]{income.png}
  \caption{Income management form with all-time and monthly totals}
  \label{fig:income}
\end{figure}

\subsection{Expense Management}

The Expenses page (Figure~\ref{fig:expenses}) follows a similar
layout to the income page: summary cards for all-time and
current-month spending and a form with name, amount, category and
date. It also lists the most recent expenses with a delete control
for each entry.

\begin{figure}[htbp]
  \centering
  \includegraphics[width=\textwidth]{expense.png}
  \caption{Expense management page with recent expenses}
  \label{fig:expenses}
\end{figure}

\subsection{Budget Management}

On the Budget page (Figure~\ref{fig:budget}) the user chooses a
category and a monthly limit. Saving a limit for a category that
already has one updates it. Each budget row shows the amount spent,
the percentage used, a progress bar, an edit and a delete control,
and an ``over budget by'' message when the limit is exceeded. The
header cards show the total monthly budget and how far spending is
over or under it.

\begin{figure}[htbp]
  \centering
  \includegraphics[width=\textwidth]{budget.png}
  \caption{Budget management with per-category usage and over-budget warnings}
  \label{fig:budget}
\end{figure}

\subsection{Savings Goals}

The Savings Goals page (Figure~\ref{fig:savings}) lets the user
create or update a goal with a name, target amount, amount already
saved and an optional deadline. Summary cards show the total
target, total saved and overall progress, and each goal has a
progress bar, the remaining amount, and Edit and Delete buttons.

\begin{figure}[htbp]
  \centering
  \includegraphics[width=\textwidth]{savings.png}
  \caption{Savings goals with progress tracking}
  \label{fig:savings}
\end{figure}

\subsection{Reports}

The Reports page (Figure~\ref{fig:reports}) generates a summary for
a chosen month and year: total income, total expenses, net savings,
a category-wise expense breakdown with percentage shares, and a
table of that month's transactions. The Export CSV button writes
the report data to a CSV file.

\begin{figure}[htbp]
  \centering
  \includegraphics[width=\textwidth]{reports.png}
  \caption{Monthly financial report with category breakdown and transaction list}
  \label{fig:reports}
\end{figure}

\subsection{Settings}

The Settings page (Figure~\ref{fig:settings}) shows the logged-in
account and provides a Change Password option.

\begin{figure}[htbp]
  \centering
  \includegraphics[width=0.7\textwidth]{settings.png}
  \caption{Settings page}
  \label{fig:settings}
\end{figure}

% ===== CHAPTER 3: RESULTS & EVALUATION =====

\chapter{Results and Evaluation}

\section{Testing Methodology}

\begin{itemize}
  \item \textbf{Functional testing:} each page was exercised
    manually using the sample dataset visible in the screenshots
    (income and expense records for October and November 2026,
    three budgets and three savings goals).
  \item \textbf{Input validation testing:} empty fields, non-numeric
    or non-positive amounts and invalid dates are rejected with a
    warning on the forms.
  \item \textbf{Consistency checks:} figures that appear on more
    than one page (totals, percentages, goal progress) were
    compared across pages.
  \item \textbf{Planned:} unit tests for the calculation functions,
    performance measurement of dashboard refresh latency, and a
    pilot with 3--5 peer testers, as defined in the
    proposal~\cite{pfmproposal}.
\end{itemize}

Table~\ref{tab:testcases} lists the manual test cases and their
outcomes.

% TODO: confirm each row against the running app before submission,
% and add more rows (budget update, goal update, CSV export, delete).
\begin{table}[htbp]
\centering
\small
\begin{tabularx}{\textwidth}{|c|l|X|X|c|}
  \hline
  \textbf{ID} & \textbf{Area} & \textbf{Input / action} &
  \textbf{Observed behaviour} & \textbf{Result} \\
  \hline
  T1 & Login & Empty username or password &
    Inline error message shown & Pass \\
  \hline
  T2 & Income form & Amount empty, non-numeric, zero or negative &
    Warning shown; nothing saved & Pass \\
  \hline
  T3 & Income form & Date not in \texttt{dd-mm-yyyy} format &
    Warning shown; nothing saved & Pass \\
  \hline
  T4 & Budget usage & Compare Budget page with sidebar card &
    146\% (Rs.\,19,000) vs 150\% (Rs.\,19,500) & Fail \\
  \hline
  T5 & Goal progress & Compare dashboard with Savings Goals page &
    16\% vs 17\% for the same goal & Fail \\
  \hline
\end{tabularx}
\caption{Manual test cases and outcomes}
\label{tab:testcases}
\end{table}

\section{Implementation Status}

Table~\ref{tab:status} compares the proposal's deliverables with
the current build.

\begin{table}[htbp]
\centering
\begin{tabular}{|l|l|}
  \hline
  \textbf{Deliverable (from proposal)} & \textbf{Status} \\
  \hline
  Login with MySQL authentication & Implemented (plain-text password) \\
  \hline
  Add / view income and expense records & Implemented \\
  \hline
  Dashboard with totals and charts & Implemented \\
  \hline
  Budget engine with per-category limits & Implemented \\
  \hline
  Savings goals with progress & Implemented \\
  \hline
  Reports and CSV export & Implemented \\
  \hline
  User registration & Not yet implemented \\
  \hline
  Recurring expenses & Not yet implemented \\
  \hline
  Password hashing & Not yet implemented \\
  \hline
  Unit tests and latency measurement & Planned \\
  \hline
  Matplotlib charts & Replaced by Tkinter canvas drawing \\
  \hline
\end{tabular}
\caption{Implementation status against the project proposal}
\label{tab:status}
\end{table}

\section{Performance Metrics}

The proposal sets a Dashboard Refresh Latency (DRL) target of
at most one second for up to 5,000 transactions. This has not yet
been measured on the current build; Table~\ref{tab:performance}
will be filled in once timing is added.

% TODO: replace "To be measured" with real numbers once you have run
% the timing and the unit tests.
\begin{table}[htbp]
\centering
\begin{tabular}{|l|r|r|}
  \hline
  \textbf{Metric} & \textbf{Target} & \textbf{Achieved} \\
  \hline
  Dashboard refresh latency (5,000 records) & $\leq 1$ s & To be measured \\
  \hline
  Calculation accuracy on fixed test set & $100\%$ & To be measured \\
  \hline
\end{tabular}
\caption{Evaluation metrics defined in the proposal}
\label{tab:performance}
\end{table}

\section{Analysis}
\label{sec:analysis}

\begin{itemize}
  \item \textbf{Objectives met:} login (without password hashing),
    transaction entry, dashboard, budgets, savings goals and
    reports are working end to end against MySQL.
  \item \textbf{Challenges:} early versions rebuilt every page on
    each navigation and showed a blank content area when frames
    were embedded incorrectly; moving to a single controller with
    stacked frames and \texttt{on\_show()} refresh hooks fixed this.
  \item \textbf{Inconsistencies found during testing:}
    \begin{itemize}
      \item The Budget page reports 146\% usage (Rs.\,19,000 of
        Rs.\,13,000) while the sidebar spending-limit card reports
        150\% (Rs.\,19,500 of Rs.\,13,000). The Budget page counts
        only categories that have a budget, whereas the sidebar
        counts all expenses, so the two figures differ by the
        Rs.\,500 spent on an unbudgeted category.
      \item Goal percentages differ between pages: the same goal
        (Rs.\,10,000 saved of Rs.\,60,000, i.e.\ 16.67\%) appears as
        16\% on the dashboard and 17\% on the Savings Goals page.
        This is consistent with one view truncating and the other
        rounding to the nearest integer.
      \item Date formats differ between forms: income and expense
        dates use \texttt{dd-mm-yyyy}, while the goal deadline uses
        \texttt{YYYY-MM-DD}.
    \end{itemize}
    The first two should be unified in a single shared calculation
    function and the third in a single date-format setting.
\end{itemize}

\section{Limitations}
\label{sec:limitations}

\begin{itemize}
  \item Passwords are currently compared as plain text and the
    database credentials are set inside the source code. This was
    accepted for the first phase and must be replaced with hashed
    passwords and external configuration before real use.
  \item There is no registration screen and no recurring-expense
    support yet.
  \item All records are entered manually; there is no import from
    bank statements or UPI apps.
  \item Several queries and calculations still live inside the GUI
    classes, so the business-logic tier is not yet fully separate.
  \item Dashboard latency and calculation accuracy have not yet been
    measured (Table~\ref{tab:performance}).
\end{itemize}

% ===== CHAPTER 4: CONCLUSION =====

\chapter{Conclusion}

\section{Summary of Work}

A working desktop Personal Finance Manager has been built with
Python, Tkinter and MySQL. It supports login, income and expense
entry, a dashboard, per-category budgets, savings goals, monthly
reports with CSV export, and a settings page. Objectives~1 to~5
are met, although the login does not yet hash passwords;
Objective~6 is partly met.

\section{Key Findings}

\begin{itemize}
  \item \textbf{Finding 1:} building the transaction-and-database
    core first, as planned in the proposal, made it straightforward
    to add budgets, goals and reports on top of the same data.
  \item \textbf{Finding 2:} a single controller with stacked frames
    and refresh hooks is a more reliable navigation pattern in
    Tkinter than recreating widgets on each page change.
  \item \textbf{Finding 3:} the same figure computed in more than
    one place drifts apart (146\% versus 150\% usage), so shared
    calculation functions are needed.
\end{itemize}

\section{Future Work and Recommendations}

\begin{enumerate}
  \item Hash passwords, move credentials to a configuration file,
    and add user registration.
  \item Add recurring expenses with due-date reminders.
  \item Move queries and calculations into a separate data-access
    and logic layer, and add unit tests for the formulas.
  \item Measure dashboard refresh latency on 5,000 records and add
    database indexes on user, date and category columns.
  \item Expose the logic layer through a Flask REST API, following
    the proposal's path from desktop MVP to web, multi-user cloud
    and mobile versions~\cite{pfmproposal}.
\end{enumerate}

\section{Final Remarks}

The project shows that a simple, auditable finance tool can be
built with standard tools and extended in stages. The next phase
is to harden security, remove duplicated calculations and
measure performance against the targets set in the proposal.

% ===== BIBLIOGRAPHY =====

\printbibliography[heading=bibintoc]

\end{document}